% ============================================================
%  Niels Treschow, "Om den menneskelige Natur i Almindelighed, især dens
%    aandelige Side."
%  Kiøbenhavn: Trykt paa Fr. Brummers Forlag hos Ernst Andreas Herman Møller,
%    Hof- og Universitetsbogtrykker, 1812.
%  TRANSCRIPTION (Danish) — THE COMPLETE WORK, in progress.
%  The first psychology written in Danish. Body: printed pp. 1–486.
%
%  STRUCTURE (from the Indhold, printed pp. XV–XVI = scan 17–18):
%    Fortale ................................................ I–VI+ (scan ~5–10)
%    Indhold ................................................ XV–XVI (scan 17–18)
%    Indledning ............................................. 1–10
%    Første Hovedstykke. Almindeligt Begreb om Mennesket,
%      især fra dets aandelige Side ......................... 11–34
%    Andet Hovedstykke. Om Forestillingsevnen
%      1ste Afdeling, om denne Evne i Almindelighed ......... 34–60
%      2den Afdeling, om Sandseligheden og umiddelbare
%        Fornemmelser ....................................... 60–90
%      3die Afdeling, om Evnen til at kalde Forestillingerne
%        tilbage, og middelbare Fornemmelser ................ 90–146
%      4de Afdeling, om Tænkekraften eller den høiere
%        Forestillingsevne .................................. 146–188
%    Tredie Hovedstykke. Om Følеevnen
%      1ste Afdeling, om Følelser i Almindelighed ........... 188–227
%      2den Afdeling, om de forskiellige Slags behagelige og
%        ubehagelige Følelser ............................... 227–269
%    Fierde Hovedstykke. Om den menneskelige Villie
%      1ste Afdeling, om Villien i Almindelighed ............ 270–293
%      2den Afdeling, om Drifter og Tilbøieligheder ......... 293–363
%      3die Afdeling, om Vaner og Færdigheder ............... 363–383
%      4de Afdeling, om Lidenskaber og Affecter ............. 384–486
%
%  Treschow numbers his §§ continuously through the whole book (1., 2., 3., …),
%  set inline at the head of each paragraph, not as display headings. The
%  Indledning runs §§ 1–…; § 8 falls at p. 16, inside the Første Hovedstykke.
%
%  WHY THIS BOOK MATTERS HERE. Its anti-Cartesian core (§§ 8–14, pp. 16–27) is
%  the nearest thing in Treschow to a doctrine of complementary modes of access:
%  the occasionalist and pre-established-harmony solutions are rejected (§ 8);
%  the two "senses" — indvortes and udvortes — give two modes of access to one
%  and the same being (§ 10, thesis at p. 20); § 11 argues from the limits of
%  those modes that the question whether bodies think is undecidable; § 12 lists
%  four senses of "Aand" including the human being as noumenon or Idea; § 13
%  answers the objection that the unity of consciousness is "kun en Form i
%  Forstanden"; and p. 27 concludes that "Mennesket er altsaa i sin individuelle
%  Grundform og som Idee eller Noümen uforgængeligt" — tying the book to the
%  Grundform of the 1810 essay and the individual Ideas of "Om Gud" (1831).
%  Companions: ../enslige-ting/, ../almindelig-logik/, ../om-gud/.
%
%  Scan: ~/bibliotek/Treschow, Niels/om-den-menneskelige-natur-1812.pdf
%        (511 scan pp.; nb.no IIIF, URN:NBN:no-nb_digibok_2008040712004;
%         public domain, "Tilgang for alle").
%  OFFSET (VERIFIED): scan page = printed + 18, and it holds across the volume:
%        p. 1 = scan 19; p. 16 = scan 34; p. 27 = scan 45; p. 486 = scan 504.
%        Spot-check when entering a new Hovedstykke.
%  TYPE: FRAKTUR. Emphasis = letterspacing (Sperrsatz) -> \emph{}; Treschow
%        letterspaces whole thesis-sentences here, not just names.
%  ORTHOGRAPHY: uses "ø"; note "Siel"/"Sielen" beside "Siæl" (both occur —
%        transcribe as printed), "Gienstand", "Bevidsthed", "Noümen" (with
%        diaeresis), "Sands" (sense), "indvortes/udvortes" (inner/outer).
%  See RESUME-NOTES.md for progress, the batch order, and the working method.
% ============================================================
\documentclass[12pt,a4paper]{article}
\usepackage[utf8]{inputenc}
\usepackage[T1]{fontenc}
\usepackage{libertinus}
\usepackage{libertinust1math}
\usepackage[danish]{babel}
\usepackage[protrusion=true,expansion=true]{microtype}
\usepackage{setspace}
\usepackage[margin=1.2in]{geometry}
\usepackage{fancyhdr}
\usepackage{hyperref}

\hypersetup{
  pdftitle={Om den menneskelige Natur i Almindelighed (1812)},
  pdfauthor={Niels Treschow},
  hidelinks
}

\pagestyle{fancy}
\fancyhf{}
\rhead{Treschow, \emph{Om den menneskelige Natur} (1812)}
\cfoot{\thepage}

\setstretch{1.3}

\newcommand{\hovedstykke}[2]{%
  \bigskip\begin{center}{\Large #1}\\[6pt]{\large \emph{#2}}\end{center}\smallskip}
\newcommand{\afdeling}[2]{%
  \medskip\begin{center}{\large #1}\\[4pt]{\emph{#2}}\end{center}\smallskip}
% Unnumbered top-level heading (Indledning), i.e. no "Hovedstykke" line above it
\newcommand{\afsnit}[1]{%
  \bigskip\begin{center}{\Large \emph{#1}}\end{center}\smallskip}

\title{\textbf{Om den menneskelige Natur i Almindelighed,}\\[4pt]
  \large især dens aandelige Side}
\author{Niels Treschow}
\date{\small Kiøbenhavn: Fr.\ Brummers Forlag, 1812}

\begin{document}
\maketitle
\thispagestyle{fancy}

\bigskip

% ############################################################
% ###  INDLEDNING  (printed pp. 1–10 = scan 19–28)         ###
% ############################################################

\afsnit{Indledning.}

% printed p. 1
\noindent
1.~\emph{Anthropologie er en Videnskab om den menneskelige Natur.} Som en
besynderlig Form blandt de uendelig mange, hvori det høieste Væsen aabenbarer
sig, eller som Naturvæsen kan Mennesket ikke kiendes uden Erfarenheds Hielp.
Denne Videnskab er følgelig ikke reen, men anvendt, og Methoden i samme mere
analytisk end synthetisk. Som egentlig Videnskab indskrænker den sig vel ikke til
at samle og ordne hvad Erfarenhed kan lære, men søger tillige at udforske de
Love, hvorefter Særsynerne rette og de Aarsager, hvoraf de reise sig, først de
nærmere, siden de fiernere, saavidt man ved en analytisk Stigen fra de
besynderlige til de almindelige Grunde i denne Række kan naae. Men, om mange Ting
end synes uforklarlige, kan de dog være alt for vigtige til med Stiltienhed at
forbigaaes: den bør der%
% printed p. 2
for alligevel optage dem, og i saadanne Tilfælde ei undsee sig ved at fremtræde
som blot Disciplin.

2.~Mennesket kan betragtes fra en dobbelt Side, nemlig en legemlig og aandelig.
Hvilkensomhelst man udelukkende giør til sit Formaal, bliver Kundskaben ensidig.
Anthropologien burde derfor, ligesom fra et høiere Sted, hvorfra begge viste sig
forenede, omfatte det hele Menneske. Men gives der da end en almindelig eller
fælles Sands, i Midten af den indvortes og udvortes, hvoraf man til denne
Betragtning kunde laane Stoffen; saa er dog hverken ved dennes saa faa, nøgne og
abstracte eller endnu mindre ved gandske rene Begreber nogen ret anskuelig
Kundskab om den menneskelige Natur at erholde: man nødes derfor til at undersøge
enhver af disse Sider for sig selv, og derefter at \emph{dele Anthropologien i
den physiske eller physiologiske og psychologiske.} Begge kan alene derved
udmærke sig fra blot Physiologie og Psychologie, at, endskiøndt i enhver af de
første en Side er den herskende, saa tager man dog heri ogsaa altid Hensyn til
den anden, og sammenligner begge; medens man i de andre Videnskaber saaledes
behandler den, som om baade Aand og Legeme var noget i sig selv, der uafhængig
kunde virke eller være til.

% printed p. 3
3.~\emph{Den psychologiske Anthropologie er deels en almindelig}, der
foredrager alt hvad den indvortes Erfarenhed om Mennesket overhovedet kan lære;
hvori der uden Tvivl forekommer Meget, som det har tilfælles enten med alle
aandelige Væsener eller flere: \emph{deels en besynderlig}, der undersøger hvad
der er Menneskets Aand eiendommeligt eller visse Slags Mennesker og Tilstande
særegent. Uagtet det nu ei er muligt saa nøie at adskille disse Dele fra
hverandre, at der jo i den første kan forekomme Meget, som egentlig vedkommer den
anden; saa gives der dog nogle saa udmærkede Forskielligheder, saavel i Tilstand,
f.~Ex.\ Alder, som i Siels og Legems oprindelige Form, f.~Ex.\ hos de tvende Kiøn,
at de fortiene en nøiere Undersøgelse, hvilken i en almindelig Afhandling ei kan
have Sted. En Videnskab om saadanne den menneskelige Naturs Særegenheder er hvad
man i Ordets strængeste Forstand kan kalde Anthropologie; hvis den forhen
bestemte Forskiel mellem denne Videnskab og Physiologie tilligemed empirisk
Sielelære skulde synes Nogen utilstrækkelig.

4.~Ethvert Væsen er som indskrænket og lidende, Naturnødvendighed, der altid er
en udvortes, og Naturlove underkastet. Da nu Sielen, ligesaa lidt som Legemet, er
% printed p. 4
gandske selvvirksom, men altid maa bestemmes af noget andet; saa ere Naturlove
ogsaa for den gieldende, dog kun for saavidt som den er lidende, d.~e.\ hvad de
Egenskaber eller Forandringer angaaer, som ei af Kraftens egen Natur og Væsen
lade sig forklare. I Anthropologien bør derfor disse Love fremsettes og anvendes.
\emph{I denne Henseende er denne Videnskab en aandelig Naturlære.}

5.~Da Philosophien tildeels derved udmærker sig fra andre Videnskaber, at det er
den indvortes Erfarenhed især, hvoraf den laaner Stoffen for sine Begreber, saa er
Sielelæren eller den psychologiske Anthropologie dermed ogsaa nærmere forenet end
nogen anden empirisk Videnskab eller Disciplin. Saa meget mere bør den derfor
beflitte sig 1) paa \emph{Fuldstændighed} i at samle de mange enslige Erfaringer,
der deels almindelig forekomme, deels sielden og ei uden anvendt Kunst eller
Flid, fremmede saavelsom egne. 2) \emph{Nøiagtighed} i at adskille saavel
Slutninger og egne Meninger derfra som andre Forestillinger, der tidt uvilkaarlig
forbinde sig med hvad man virkelig erfarer. Dernæst at bemærke hver enslig
Omstændighed, som til at bedømme det Skedte kan være nødvendig, og i Udtrykket
være, saa meget mueligt er, bestemt. 3) \emph{Orden}, hvortil udfordres
% printed p. 5
en Regel, hvorefter man ei alene samler, hvad der indbefattes under et
almindeligt Begreb, men ogsaa hvad der kan udledes af samme Aarsager eller
Grunde. Uden saadan Regel blive Begreberne om de forskiellige Sindets Evner og
Inddelingen deraf ganske forvirrede, som man i det Galske System seer et
mærkværdigt Exempel paa. 4) \emph{Grundighed}, der i det mindste veed at angive
de nærmeste Grunde. Disse ligge deels i de Hovedevner, hvorfra man pleier at gaae
ud, deels i de almindelige Naturlove, deels i Legemets Organisation eller det
aandelige Væsens Analogie med samme. \emph{Overfladig} bliver Anthropologien
derimod, naar den af faa, tidt misforstaaede, Erfaringer uddrager almindelige
Slutninger, uden at tage deres Analogie i Betragtning eller have tilbørligt
Hensyn til de høiere og høieste noksom afgiorte Grundsætninger, hvorpaa al
Videnskab saavel om Aand som Legeme beroer.

6.~\emph{Kundskab om Mennesket, især dets aandelige Side, er baade i sig selv af
største Vigtighed, og uundværlig: 1) som Middel til at opnaae vor Bestemmelse.}
Er dette Lyksalighed, saa maae vi nødvendig kiende vore Følelsers Beskaffenhed og
Grunde. Ubekiendtskab med de Anlæg, Naturen har for%
% printed p. 6
synet os med, for at afvende Smerter og uden Nag at nyde, er aabenbar en
Hovedaarsag, hvorfor Mange forfeile dette Maal. Kiendte man bedre de mange Kilder
til rolige og altid stigende Fornøielser, som Sandheds Kundskab og den ideale
Verdens Betragtning aabner, vilde man da vel ved en umaadelig Nydelse af hvad der
kildrer den grove Sandselighed forbittre det hele følgende Liv? Er derimod
\emph{Sædelighed} det høieste Gode; saa ligger Roden dertil i Fornuften og den
fri Villie: hvilken til alle Tider har været en af de vigtigste Formaale for
Philosophers og især for Psychologers Eftertanke. Men Mennesket er tillige
sandseligt. Det har Tilbøieligheder, som det ved Fornuften skal regiere; hvilket
er umuligt uden at kiende dem. Moralphilosophie og Sielelære ere derfor saa nøie
forenede, at man i den første ingen Fremgang kan giøre uden dyb Indsigt i den
anden. \emph{Men i flere Videnskaber er 2) saadan Kundskab ligeledes nødvendig,
saasom:} a) \emph{i den practiske Religionslære.} Menneskekundskab har man altid
holdt for en Egenskab, der især bør findes hos Folkelæreren, som til Lasters,
Daarligheders og Fordommes Udryddelse, til sand Oplysnings og Dyds Udbredelse,
til Trøst, Raad og Afskrækkelse ligesaa vel bør vide at overtale, som at
overbevise, og til den Ende maa kiende de subjective Grun%
% printed p. 7
de, der hos Folket ere de mest virksomme. Men er denne Kundskab blot empirisk,
bliver den deels ensidig, og forleder til hyppige Feiltagelser, deels er den af
Mangel paa Principier i mange Tilfælde utilstrækkelig. b) \emph{De vigtigste
Spørgsmaale, der i Pædagogiken forekomme, kan man uden Sielelærens Hielp ikke
besvare.} Dertil hører det om Ordenen og Maaden, paa hvilken Børns Begreber bør
udvikles, og nye Kundskaber dem bibringes, om Hiertets Dannelse, om Belønninger
og Straffe, med mange flere. c) \emph{De skiønne Videnskabers Theorie} ville nogle
Nyere vel blot have grundet paa Ideer eller rene Begreber og et høiere Slags
Geniet særegen Anskuelse. Men er det end ganske rigtigt, at Skiønhedens Kierne,
saa at sige, er en Idee; saa maa denne dog fremstilles i en for os behagelig Form,
ei efter en besynderlig Modes Forskrift, men efter Beskaffenheden af vor Anskuelse
i Almindelighed: hvilken man hverken af Exempler eller rene Begreber alene, men af
den menneskelige Natur maa lære. Mangel af denne Kundskab har man endog i de
seneste Tider en Mængde Uhyrer at takke og en ved i sig selv rigtige
Grundsætningers Misforstaaelse gandske fordærvet Smag. d) \emph{Der gives neppe
tvende Videnskaber, der mere trænge til hverandre end Sielelæren og Lægekunsten.}
Det er derfor ogsaa
% printed p. 8
enten virkelige Læger eller Mænd, udrustede med samme Slags theoretiske
Kundskaber, som i vore Dage have leveret de mest interessante Bidrag til
Anthropologiens, endog den psychologiskes, Uddannelse. Men ligesaa nødvendig er
igien denne philosophiske Videnskab i Lægekunsten selv. Grunden til mange
Sygdomme ligger mere i visse herskende Forestillinger og deres Forbindelse end i
nogen Legemets, d.~e.\ den grovere Organismes, Uorden. Ofte har en glad Tidende,
Adspredelse, Frygt eller Haab overvundet haardnakkede Sygdomme, der have trodset
de kraftigste Lægemidler. Den lykkeligste Læge er gemeenlig den, den bedst
forstaaer at omgaaes den Syge, der med dyb Menneskekundskab forener Taalmodighed i
at studere hans Tilbøieligheder og Charakter. e) \emph{Hvorvidt Sielelæren paa den
øvrige Naturlære kan anvendes}, synes ei alene tvivlsomt; men Nogles uheldige
Forsøg kunde endog afskrække derfra. Er det alligevel vist, at Kundskab om
Legemet er skikket til at udbrede Lys over den aandelige Natur, den udvortes
Erfarenhed over den indvortes; saa kan det neppe feile, at denne igien maa kunne
bevise hiin samme Tieneste, om end nærværende Tidsalders herskende Tendens gaaer
heller den modsatte Vei, og derfor ei er meget tilbøielig til at agte paa denne
Nytte.
% printed p. 9
f) \emph{Hvad Lovkyndigheden angaaer}, saa er først dette hele Studium gandske
grundet paa Moralphilosophie, og denne igien paa Sielelære. Dernæst vil neppe
nogen nægte, at jo den Fuldkommenhed, som man indtil de senere Tider i Criminel-
og Procesretten saa høiligen har savnet, har været Mangel af denne Kundskab at
tilskrive. Lovgiverens Hensigt bør det ikke være blot at hævne eller straffe, men
at forekomme Forbrydelser og afskiære Leiligheden til voldsomme Udbrud af
Lidenskaber, hvoraf de fleste Misgierninger reise sig. Er vel dette mueligt uden
dyb Indsigt i den menneskelige Natur saavel i det Hele, som i alle de Stykker,
hvorpaa Forskiellen mellem Mennesker beroer? Dommeren skal opdage Sandhed, lokke
den Skyldige til Bekiendelse, tvinge List og Underfundighed til at aflægge Masken,
eller søge, uagtet al List og Forstillelse, at blive skiulte Tanker og Handlinger
vaer. Ukyndighed har taget sin Tilflugt til Pinebænke. Hvem har afskaffet dem uden
Philosophien, som her beskæftiger sig med Undersøgelsen af det menneskelige
Hierte? g) \emph{Ustridig er Menneskekundskab i Statskunsten det allervigtigste.}
Af Uvidenhed om samme har man giort List og Klogskab til ensbetydende. Overbeviist
om, at der ingen anden Drifter gives end egennyttige, har Stats%
% printed p. 10
manden kun ved Rænker og Forblindelse troet at kunne naae sit Maal. Hvad der i
hiin Sætning er rigtigt, nemlig at disse Drifter hos de Fleste ere de herskende,
har man overdrevet, og ved at omgaaes andre Mennesker efter denne Grundsætning,
fornedret dem, saaret deres moralske Følelse og givet dem et Slags Ret til at
behandle Næsten igien paa samme Maade.

\begin{center}\rule{0.22\linewidth}{0.4pt}\end{center}

% ############################################################
% ###  FØRSTE HOVEDSTYKKE  (printed pp. 11–34 = scan 29–52) ###
% ############################################################

\hovedstykke{Første Hovedstykke.}{Almindeligt Begreb om Mennesket, især fra dets
aandelige Side.}

% [Trykte s. 11–15 = scan 29–33 — endnu ikke transskriberet.
%  §§ 3–7 falde her; § 8 begynder s. 16.]

% printed p. 16
\noindent
[\dots] forklare hvad der under Forudsætning af en ikke blot formal, men
væsentlig Forskiel mellem Siel og Legeme aldeles ikke lader sig forklare. Det
Cartesianske Understøttelses eller Malebrancheske Leilighedsaarsagernes System,
der, isteden for virkende Naturaarsager, udleder den tilsyneladende
Vexelvirkning mellem begge Væsener af Guds Almagt, som af Sielens Villie og
Legemets Bevægelser kun tager Anledning til selv at virke enten paa den første
eller det andet: denne Mening, foruden at den efterlader Sagen i samme Mørke som
den forrige, indvikler sig derhos i aabenbare Modsigelser, og sætter Dualismens
Urigtighed i et endnu klarere Lys. \emph{Leibniz's forudbestemte Harmonie}
saavel mellem Siel og Legeme som mellem andre Substanser, der ved indvortes
Kraft alene selvstændig skal frembringe alle deres Forandringer, uden mindste
Stød af fremmede Ting, kan vel i Sammenligning med hans øvrige Lærdomme maaskee
saaledes udtydes, at man fra en vis Side maa bifalde den; men som en blot physisk
eller anthropologisk Hypothese mangler den alligevel saavel indvortes som
udvortes Sandsynlighed.

8.~Det Spørgsmaal om Siælens Sæde har forekommet Nogle ei alene fra vor
Standpunct, men ogsaa i sig selv umuligt at besvare, og derfor urimeligt at
opkaste, da Sielen efter deres Mening som en Gienstand enten for den høiere
Tænken alene eller for den indvortes Sands ei er bunden til noget Sted. Men,
skiøndt
% printed p. 17
det ei kan nægtes, at Gienstande for begge Slags Erfarenhed saavel indbyrdes som
endnu mere fra de inteligible eller virkelige Ting ere meget forskiellige; saa
maa Enheden deri dog nødvendig overalt aabenbare sig, og Sielens eller det
egentlige Menneskes altsaa just som en synlig Punct, fra hvilken det synes, at al
Virksomhed gaaer ud. Har man derfor end hidindtil ikke fundet samme, enten i
Kronkirtelen, i den saa kaldte Varols Bro, i Hiernehulerne, i den forlængede Marv
eller paa noget andet Sted; saa er den dog uden Tvivl et Sted at finde, og
Opdagelsen deraf er, som et empirisk Beviis eller udvortes Symbol paa den Enhed,
som Bevidstheden eller den indvortes Sands langt tydeligere fremstiller, vel
ingenlunde at foragte.

9.~\emph{Af langt større Vigtighed er det at vide, hvorvidt den Kundskab, vi om
Legemets Organisation eller det udvortes Menneske kan have, nytter til at oplyse
vor aandelige eller det indvortes Menneskes Natur, eller om tvertimod Lovene for
begges Virkninger ere saa forskiellige, at man ei kan overføre dem fra det ene
til det andet Fag.} At nogen caussal Forbindelse mellem dem ei kan tænkes, er nu
vel af det Foregaaende noksom aabenbart. Men, uagtet al Forskiellighed, tillader
Væsenets Identitet dog ingen Tvivl om en vis Analogie. Sproget selv viser
tydelig, at man her længe siden
% printed p. 18
har anet den; da man ellers ei saa hyppig vilde have laant Ord og Udtryk, især
af legemlige Ting, naar man manglede noget egentligt for at betegne Sielens
Egenskaber og Virkninger, ligesom der ogsaa gives Exempler paa det Modsatte.
Saadan Ligheds Forudsettelse i de Tilfælde, hvor Menneskets ene Side maa være
mere bekiendt end den anden, kan ei alene føre til nye Loves Opdagelse og de
besynderlige Videnskabers Berigelse, der have enten Siel eller Legeme til
Formaal, men endog til at bringe os den Punct nærmere, hvori begge eller deres
Hovedyttringer, Bevægelse og Fornemmelse, støde sammen: hvilket uden Tvivl er al
Videns, som saadans, høieste Maal. De Indvendinger, som Dualisterne herimod have
giort, bevise kun, hvor ufuldkommen vor Kundskab om denne Sag hidindtil er
bleven, uagtet alle Bestræbelser for at brække de Skranker ned, der holde vor
Videlyst tilbage. Men paa den anden Side viser baade den Fremgang, som den
psychologiske Anthropologie virkelig har Naturforskeres Flid at takke, og den
Fornøielse, som Opdagelsen af enhver skiult Analogie medfører, endnu langt
tydeligere, at man om et lykkeligt Udfald af denne Flid ei bør mistvivle.

10.~\emph{Den sædvanlige Deling af Væsener i Legemer og Aander beroer efter det
dualistiske System paa følgende Grunde. a) Vi have to Slags Sands, indvortes og
udvortes.}
% printed p. 19
\emph{Ved hiin erfare vi umiddelbar en Virksomhed, Attraa, Stræben, Lyst og
Ulyst, der alle gaae ud paa visse Forestillinger; ved denne mærke vi Lys, Farve,
Figur, Bevægelse, Lyd, Toner, Vellugt, Surt, Sødt, Haardt, Blødt o.~s.~v.} Nu ere
disse Fornemmelser vel indbyrdes næsten ligesaa forskiellige som fra det første
Slags, hvori der ogsaa er megen Forskiellighed. Thi hvad Sammenligning kan man
vel giøre mellem Farven og Lugten af en Rose, mellem Formen og Lyden af en
Fiolin? Ikke desmindre lærer Erfarenhed, at de alle komme frem af de samme
Gienstande. Ingen kan tvivle paa, at det jo er den samme Rose, der lugter, som
den, der indtager Synet ved sin Farve. Dette overbeviser os om, at der ere
Legemer til, og at alle disse Fornemmelser, hvor forskiellige de end maae være,
dog ikke reise sig af Gienstandenes egen indvortes Mangfoldighed, men Sandsens,
ved hvilken vi fornemme den. Derimod have vi af Erfarenhed ingen Grund til at
troe, at disse Legemer tillige have Forestillinger, tænke, føle, ville, eller paa
den anden Side, at vi, der saaledes afficeres eller handle, selv have nogen
Figur, Farve, Lugt eller Smag. Men at alt hvad vi ved den udvortes Sands
fornemme tilhører og kommer fra et eneste Subject, at Tanker, Følelser,
Begieringer have deres Udspring af det samme, ei af forskiellige, derom blive vi
af vor Selvbevidsthed fuldkommen overbeviste. Naar derfor nogle gamle
Philosopher, ja ogsaa nogle nyere
% printed p. 20
have lært, at vi have flere Siele, to eller tre; saa maa de enten ikke have agtet
derpaa, eller, hvad der synes rimeligere, dog antaget en saadan Forbindelse
mellem dem, hvorved de ikke destomindre maatte betragtes som en eneste.
\emph{Det der fornemmes med den indvortes Sands, er altsaa Aand, Legeme hvad vi
alene ved den udvortes kan fornemme.}

b) Materien, hvoraf Legemverdenen bestaaer, virker vel paa vore Sandser, men, som
det synes, ved sin blotte Tilværelse, ei ved nogen Stræben, Drift og
Selvbestemmelse. De fleste Legemer ere saaledes beskafne, at de hverken af sig
selv kan bevæge sig, eller, naar de af andre settes i Bevægelse, gaae videre end
disse drive dem. Er der end nogle, som de organiske, der yttre Tegn til et Slags
Selvvirksomhed, saa kan Livet dog igien forlade dem, og de falde da i den forrige
døde Tilstand tilbage. Livet er altsaa en usynlig Kraft; thi Døden formindsker
ikke Massen eller Materien selv, ja, Formen bliver endog i kortere eller længere
Tid uforandret, efterat dette Liv er borte. Saaledes faae vi Begreb om Aand som
Noget, hvis Væsen er Selvvirksomhed, og vi ere tilbøielige til at troe, den
findes overalt, hvor vi ei af udvortes og mechaniske Aarsager kan forklare de
Forandringer, som i Legemerne foregaae. Derfor ansaae de Gamle den hele Natur for
besielet. Det hylozoistiske System er for den almindelige Forstand det
% printed p. 21
fatteligste, og i Philosophien selv maa man af høiere Grunde vel ogsaa omsider,
skiøndt i en finere Form, optage det igien. \emph{Selvvirksomhed og Mangel paa
samme er følgelig det andet væsentlige Kiendemærke, der bestemmer Forskiellen
mellem Legeme og Aand.}

11.~\emph{Men ere disse Kiendemærker vel i den strængeste Forstand saa rigtige,
at ingen Overgang er mulig fra det ene Slags Væsener til det andet? For det
første kan om Tingenes negative og privative Beskaffenheder Erfarenhed intet
lære.} Om Legemer kan tænke, Aander antage nogen Skikkelse maa Fornuft afgiøre.
Deraf, at vi intet saadant erfare, kan man saa meget mindre slutte det Modsatte,
som det er klart, at, om Legemer end virkelig tænkte eller følte, Aander derimod
ved en vis Gestalt gave sig tilkiende; saa kunde vi dog hverken ved den indvortes
Sands fornemme det første eller ved den udvortes det andet, og at følgelig Sagen
paa denne Maade nødvendig maatte blive skiult. Vel kan man sige, at de
Egenskaber, hvorved begge Slags Væsener udmærke sig, ere gandske stridige, og
følgelig i samme Subject ei tilsammen kan bestaae. Men lader sig ei samme
Indvending giøre mod Aand og Legeme selv, hvori man ved enhver Art af Erfarenhed
ligeledes opdager mange tilsyne stridige Ting, f.~Ex.\ i den første Sand%

\end{document}
